\chapter{Introduction and Motivation}%
\label{chap:Introduction}%


Autonomous vehicles are rapidly becoming a reality. However, their perception systems still face significant challenges 
in complex and unforeseen scenarios, often leading to disengagements, where human intervention via teleoperation is required. 
In such critical moments, the vehicle must rely on an accurate and robust perception system to ensure safe navigation. 

Most existing perception systems, responsible for detecting and tracking objects in the vehicle's surroundings, rely heavily
on learning-based solutions. While these methods have demonstrated impressive performance in structured and controlled environments, 
they often struggle to generalize in novel and untrained scenarios. This limitation is particularly problematic in safety-critical
situations, where the vehicle must make quick decisions based on the perceived environment.

This project aims to develop a deterministic, stereo camera-based perception system for autonomous vehicles, capable of providing
a reliable and robust understanding of the vehicle's surroundings, even in highly dynamic and unforeseen scenarios. The system
will take as input two images from a stereo camera mounted on the vehicle and output a list of detected objects in the scene, along
with their criticality scores, forming a situational awareness module to assess potential hazards. 
Additionaly, this deterministic approach could run in parallel with existing learning-based perception systems, providing a
an additional layer of redundancy and security in critical situations where the learning-based methods may fail.

To achieve accurate 6D motion estimation (position + velocity), the system will implement a deterministic algorithm 
inspired by the Daimler 6D Vision system. This approach uses Kalman Filters to estimate a 6D state vector for selected feature
points in the scene, which will then be clustered to identify moving objects. Unlike learning-based methods, this approach
does not require extensive training data and can generalize well in new unseen scenarios.

Given the real-time requirements of autonomous vehicles, the system will be implemented in C++ using the Robot Operating System (ROS)
framework. It will also incorporate GPU hardware acceleration using CUDA to speed up computationally intensive tasks, such as
feature point detection, clustering or optical flow computation.  

To ensure scalability and modularity, the entire system will be containerized using Docker, facilitating deployment across different
hardware configurations. Finally, the developed system will be deployed and evaluated on the EDGAR research vehicle. Here, 
the system's performance will be compared against the state-of-the-art learning-based perception systems currently used in the vehicle.

